\chapter{Äquivalente Fassungen des Auswahlprinzips}
\label{b09guide:ch:criteria}

Die folgenden Fassungen beschreiben denselben Existenzgrundsatz.
Ihre Äquivalenz ist über ZF gemeint, also vor der Hinzunahme
des Auswahlaxioms. \cite{hrbacek,spronk}

\section{Auswahl aus beliebigen Mengenfamilien}

\begin{theorem}[Familienform mit Auswahlfunktion]
\label{b09guide:thm:family}
Unter AC besitzt jede Mengenfamilie $\mathcal A$ ohne leeres Glied
eine Auswahlfunktion:
\[
 \varnothing\notin\mathcal A
 \quad\Longrightarrow\quad
 \exists c\colon\mathcal A\longrightarrow\bigcup\mathcal A\quad
 \forall A\in\mathcal A\;(c(A)\in A).
\]
Die Familienglieder dürfen sich beliebig überlappen.
Die allgemeine Aussage ist äquivalent zu AC.
\end{theorem}

Eine Auswahlfunktion verlangt genau einen \emph{gewählten Wert}
je Familienglied. Sie verlangt nicht, daß die Menge aller gewählten
Werte jedes überlappende Familienglied in nur einem Element trifft.

\begin{theorem}[Indizierte Auswahlform]
\label{b09guide:thm:indexed}
Unter AC gilt für jede indizierte Mengenfamilie $(A_i)_{i\in I}$:
\[
 \begin{gathered}
 \forall i\in I\;(A_i\ne\varnothing)\\
 \Longrightarrow\
 \exists f\colon I\longrightarrow\bigcup_{i\in I}A_i\quad
 \forall i\in I\;(f(i)\in A_i).
 \end{gathered}
\]
Dabei ist $I$ eine beliebige Menge; die $A_i$ müssen nicht verschieden
sein. Die allgemeine Aussage ist äquivalent zu AC.
\end{theorem}

\section{Auswahl auf der Potenzmenge}

\begin{theorem}[Einheitliche Auswahl innerhalb einer festen Menge]
\label{b09guide:thm:powerset}
Unter AC besitzt jede Menge $U$ eine Auswahlfunktion auf allen
ihren nichtleeren Teilmengen:
\[
 \ChoiceSet(U)\ne\varnothing.
\]
Die Aussage für alle Mengen $U$ ist äquivalent zu AC.
\end{theorem}

Die Quantoren lauten $\forall U\,\exists c\in\ChoiceSet(U)$.
Es wird keine einzige Klassenfunktion auf allen nichtleeren Mengen
behauptet; ein Prinzip der globalen Auswahl gehört nicht zu
diesem Leitfaden. Auch eine eindeutige oder kanonische Auswahlregel
ist damit nicht festgelegt.

\section{Nichtleere Produkte}

\begin{theorem}[Produktform und Kriterium für das leere Produkt]
\label{b09guide:thm:product}
Unter AC gelten für jede indizierte Mengenfamilie $(A_i)_{i\in I}$
die Äquivalenzen
\[
 \begin{aligned}
 \prod_{i\in I}A_i\ne\varnothing
   &\quad\Longleftrightarrow\quad
        \forall i\in I\;(A_i\ne\varnothing),\\
 \prod_{i\in I}A_i=\varnothing
   &\quad\Longleftrightarrow\quad
        \exists i\in I\;(A_i=\varnothing).
 \end{aligned}
\]
Die Aussage, daß jedes Produkt nichtleerer Mengen nichtleer ist,
ist äquivalent zu AC.
\end{theorem}

Der Ausdruck \emph{leeres Produkt} kann die leere Indexmenge oder
eine leere Produktmenge meinen. Hier bleibt beides getrennt:
Für $I=\varnothing$ ist die Produktmenge $\{\varnothing\}$ und
damit nicht leer. Für nichtleere Indexmengen kann das Scheitern
einer gemeinsamen Auswahl ohne AC nicht allgemein auf einen
leeren Faktor zurückgeführt werden.
